\documentclass[11pt,a4paper]{article}
\usepackage[margin=2.2cm]{geometry}
\usepackage{microtype}
\usepackage{booktabs}
\usepackage{xcolor}
\usepackage{pgfplots}
\usepackage{pgfplotstable}
\usepackage[hidelinks]{hyperref}
\pgfplotsset{compat=1.18}

\definecolor{pos}{HTML}{1B7F3B}
\definecolor{neg}{HTML}{B3261E}

% Generated by `java -cp out Backtest`
\newcommand{\NPicks}{50}
\newcommand{\MeanRet}{+17.4\%}
\newcommand{\MedianRet}{+11.4\%}
\newcommand{\SdRet}{30.7\%}
\newcommand{\Winners}{35}
\newcommand{\Losers}{15}
\newcommand{\HitRate}{70\%}
\newcommand{\BestSym}{INTC}
\newcommand{\BestRet}{+92.9\%}
\newcommand{\WorstSym}{UNH}
\newcommand{\WorstRet}{-42.6\%}
\newcommand{\TopTenRet}{+19.6\%}
\newcommand{\RestRet}{+16.9\%}


% Fill in the test window here once, it is used throughout the report.
\newcommand{\Window}{the entry and exit dates of the test window}

\title{A Rank-Sum Stock Screener on S\&P~500 Fundamentals}
\author{stock-ranker}
\date{\today}

\begin{document}
\maketitle

\begin{abstract}
A Java program pulls quarterly financial statements, company overviews, news
sentiment and earnings surprises from the Alpha Vantage API for a list of
S\&P~500 tickers, ranks every company on 50 metrics, and combines the ranks
into one weighted score. The \NPicks{} highest-scoring companies returned
\MeanRet{} on average (median \MedianRet) between \Window{};
\Winners{} of \NPicks{} rose and \Losers{} fell.
\end{abstract}

\section{Method}

For every ticker the program collects:
\begin{itemize}\itemsep0pt
  \item \textbf{Income statement, balance sheet, cash flow} -- 29 line items,
        each averaged over the last 12 quarters;
  \item \textbf{Company overview} -- 15 valuation, profitability, risk and
        analyst-rating fields (P/E, P/B, P/S, dividend yield, ROE, beta, \dots);
  \item \textbf{Derived ratios} -- gross and net margin, current ratio, debt-to-equity;
  \item \textbf{News sentiment} -- mean score of up to 25 articles from the last six months;
  \item \textbf{Earnings surprise} -- mean EPS surprise over the last 8 quarters.
\end{itemize}

For each metric $m$ with weight $w_m$, the $N$ tickers that report it are
sorted best-first (descending, or ascending for ``lower is better'' metrics
such as P/E, beta, liabilities or analyst sell ratings). The company in
position $i$ (starting at $0$) receives $w_m (N - i)$ points, and a company's
score is the sum over all metrics:
\[
  S_c = \sum_{m} w_m \bigl(N - \mathrm{rank}_m(c)\bigr).
\]
Weights range from 1 to 6. News sentiment and earnings surprise carry the
largest weight (6); net income, R\&D, earnings growth, beta, strong-buy ratings,
investing cash flow and long-term debt issuance carry 3.

\section{Results}

Table~\ref{tab:picks} lists the \NPicks{} companies in the order the screener
ranked them, with the entry and exit prices and the simple return.
Figure~\ref{fig:returns} shows the same returns sorted from best to worst.

\begin{center}
\begin{tabular}{lr}
  \toprule
  Mean return (equal-weight portfolio) & \MeanRet \\
  Median return                        & \MedianRet \\
  Standard deviation                   & \SdRet \\
  Winners / losers                     & \Winners{} / \Losers{} (\HitRate{} hit rate) \\
  Best pick                            & \BestSym{} (\BestRet) \\
  Worst pick                           & \WorstSym{} (\WorstRet) \\
  Mean of ranks 1--10 vs.\ 11--\NPicks & \TopTenRet{} vs.\ \RestRet \\
  \bottomrule
\end{tabular}
\end{center}

\begin{figure}[!ht]
\centering
\begin{tikzpicture}
\pgfplotstableread[col sep=comma]{returns.csv}\data
\pgfplotstablesort[sort key=ret,sort cmp=float >]\sorted{\data}
\begin{axis}[
    width=\textwidth, height=6.2cm,
    ybar, bar width=4.2pt,
    xtick={0,1,...,49}, % one tick per pick (0 .. NPicks-1) xmin=-0.7, xmax=\NPicks-0.3, xticklabels from table={\sorted}{symbol},
    x tick label style={rotate=90, font=\tiny},
    ylabel={Return (\%)}, ymajorgrids, grid style={gray!25},
    axis x line*=bottom, axis y line*=left,
]
\addplot[fill=pos, draw=none, bar shift=0pt, unbounded coords=jump]
  table[x expr=\coordindex, y expr={\thisrow{ret} >= 0 ? \thisrow{ret} : nan}]{\sorted};
\addplot[fill=neg, draw=none, bar shift=0pt, unbounded coords=jump]
  table[x expr=\coordindex, y expr={\thisrow{ret} < 0 ? \thisrow{ret} : nan}]{\sorted};
\draw[gray] (axis cs:-1,0) -- (axis cs:\NPicks,0);
\end{axis}
\end{tikzpicture}
\caption{Return of each pick, sorted.}
\label{fig:returns}
\end{figure}

\begin{table}[!ht]
\centering\small
\caption{The \NPicks{} top-ranked companies. Return in \%.}
\label{tab:picks}
\begin{tabular}{rlrrr@{\hspace{2em}}rlrrr}
\toprule
\# & Ticker & Entry & Exit & Ret. & \# & Ticker & Entry & Exit & Ret. \\
\midrule
1 & F & 9.71 & 13.22 & \textcolor{pos}{+36.1} & 26 & DELL & 91.46 & 153.26 & \textcolor{pos}{+67.6} \\
2 & MRK & 112.31 & 84.20 & \textcolor{neg}{-25.0} & 27 & CSCO & 45.83 & 72.08 & \textcolor{pos}{+57.3} \\
3 & XOM & 114.77 & 114.83 & \textcolor{pos}{+0.1} & 28 & APD & 279.21 & 237.63 & \textcolor{neg}{-14.9} \\
4 & TMO & 598.03 & 568.27 & \textcolor{neg}{-5.0} & 29 & CMCSA & 39.35 & 27.63 & \textcolor{neg}{-29.8} \\
5 & AMZN & 161.02 & 250.44 & \textcolor{pos}{+55.5} & 30 & ABT & 109.76 & 124.33 & \textcolor{pos}{+13.3} \\
6 & JNJ & 161.25 & 186.15 & \textcolor{pos}{+15.4} & 31 & BMY & 46.85 & 46.04 & \textcolor{neg}{-1.7} \\
7 & JPM & 194.90 & 312.27 & \textcolor{pos}{+60.2} & 32 & AIG & 72.14 & 77.82 & \textcolor{pos}{+7.9} \\
8 & WMT & 67.59 & 102.65 & \textcolor{pos}{+51.9} & 33 & HD & 348.64 & 374.83 & \textcolor{pos}{+7.5} \\
9 & TXN & 181.19 & 163.06 & \textcolor{neg}{-10.0} & 34 & INTC & 19.71 & 38.03 & \textcolor{pos}{+92.9} \\
10 & NOC & 492.63 & 573.48 & \textcolor{pos}{+16.4} & 35 & MDT & 82.04 & 89.84 & \textcolor{pos}{+9.5} \\
11 & LLY & 775.26 & 951.52 & \textcolor{pos}{+22.7} & 36 & PEP & 172.62 & 142.16 & \textcolor{neg}{-17.6} \\
12 & SO & 86.57 & 92.10 & \textcolor{pos}{+6.4} & 37 & CB & 273.41 & 284.09 & \textcolor{pos}{+3.9} \\
13 & VZ & 39.80 & 39.66 & \textcolor{neg}{-0.4} & 38 & PCAR & 95.17 & 99.99 & \textcolor{pos}{+5.1} \\
14 & AMGN & 325.22 & 319.32 & \textcolor{neg}{-1.8} & 39 & CRM & 262.91 & 254.03 & \textcolor{neg}{-3.4} \\
15 & QCOM & 157.95 & 179.10 & \textcolor{pos}{+13.4} & 40 & ALL & 191.60 & 195.70 & \textcolor{pos}{+2.1} \\
16 & GS & 459.02 & 793.51 & \textcolor{pos}{+72.9} & 41 & T & 26.24 & 24.65 & \textcolor{neg}{-6.1} \\
17 & UNH & 569.96 & 327.03 & \textcolor{neg}{-42.6} & 42 & ACN & 390.22 & 246.73 & \textcolor{neg}{-36.8} \\
18 & CAT & 316.80 & 566.89 & \textcolor{pos}{+78.9} & 43 & BAC & 46.01 & 52.82 & \textcolor{pos}{+14.8} \\
19 & IBM & 183.31 & 303.01 & \textcolor{pos}{+65.3} & 44 & CVX & 157.23 & 153.65 & \textcolor{neg}{-2.3} \\
20 & MSFT & 395.15 & 508.94 & \textcolor{pos}{+28.8} & 45 & C & 83.94 & 101.98 & \textcolor{pos}{+21.5} \\
21 & T & 19.40 & 24.65 & \textcolor{pos}{+27.1} & 46 & GOOG & 187.13 & 283.37 & \textcolor{pos}{+51.4} \\
22 & AXP & 233.54 & 365.41 & \textcolor{pos}{+56.5} & 47 & PG & 166.49 & 147.19 & \textcolor{neg}{-11.6} \\
23 & BLK & 854.83 & 1079.56 & \textcolor{pos}{+26.3} & 48 & SCHW & 82.01 & 93.84 & \textcolor{pos}{+14.4} \\
24 & DE & 349.14 & 474.44 & \textcolor{pos}{+35.9} & 49 & EIX & 52.42 & 56.95 & \textcolor{pos}{+8.6} \\
25 & AAPL & 213.31 & 270.56 & \textcolor{pos}{+26.8} & 50 & COF & 209.82 & 222.85 & \textcolor{pos}{+6.2} \\
\bottomrule
\end{tabular}


\medskip
\footnotesize AT\&T (T) appears twice (ranks 21 and 41) with different entry
prices, as in the original output; both rows are kept in the statistics.
\end{table}

\section{Limitations}
\begin{itemize}\itemsep0pt
  \item One period, no benchmark: the result should be compared with the
        S\&P~500 over the same window before drawing conclusions.
  \item Look-ahead bias: Alpha Vantage returns \emph{current} fundamentals, so
        a ranking run after the entry date partly uses information that was
        not available at entry.
  \item The ticker list is an older S\&P~500 snapshot; companies that were
        delisted or renamed simply receive no data (survivorship bias).
  \item Size-type metrics (revenue, total assets, market cap) favour large
        companies, so the screen leans towards mega-caps.
\end{itemize}

\noindent\emph{This is a student project, not investment advice.}

\end{document}
